\documentclass[11pt,letterpaper]{article}
\usepackage[T1]{fontenc}
\usepackage{lmodern,microtype}
\usepackage[margin=.68in,headheight=15pt]{geometry}
\usepackage{amsmath,amssymb,mathtools}
\usepackage{xcolor,booktabs,array,enumitem,fancyhdr,hyperref}

\definecolor{ink}{HTML}{222834}
\definecolor{indigo}{HTML}{25355D}
\definecolor{muted}{HTML}{505A6B}
\definecolor{keyink}{HTML}{205D52}
\color{ink}
\hypersetup{colorlinks=true,linkcolor=indigo,urlcolor=indigo,
  pdftitle={Conditional Probability and Bayes' Rule: 50-minute Quiz},
  pdfauthor={Chance, Data, and Decisions}}
\setlength{\parindent}{0pt}
\setlength{\parskip}{3pt}
\setlength{\abovedisplayskip}{6pt}
\setlength{\belowdisplayskip}{6pt}
\clubpenalty=10000
\widowpenalty=10000
\setlist[enumerate]{label=(\alph*),leftmargin=1.7em,
  itemsep=3pt,topsep=4pt,parsep=2pt}
\setlist[itemize]{leftmargin=1.5em,itemsep=1pt,topsep=3pt,parsep=0pt}

\newcommand{\Prob}{\mathbb P}
\newcommand{\edition}{\ifanswers Instructor key\else Student quiz\fi}
\newcommand{\partpts}[1]{\unskip\nobreak\hspace{.3em}%
  \mbox{\small\color{muted}(#1 points)}}
\newcommand{\credit}[1]{\par{\footnotesize\color{muted}\textit{Credit:} #1}}
\newcommand{\problem}[4]{\par\addvspace{8pt}%
  {\large\bfseries\color{indigo}#1.\ #2}\hfill
  {\small\color{muted}\textbf{#3 points} $\cdot$ about #4 minutes}\par\nobreak}
\NewDocumentEnvironment{solution}{O{.45in}+b}{%
  \ifanswers
    \par\begingroup\small\color{keyink}\textbf{Solution.} #2\par\endgroup
  \else
    \par\nobreak\vspace*{#1}
  \fi
}{}

\pagestyle{fancy}
\fancyhf{}
\fancyhead[L]{\small\color{muted}Conditional probability quiz}
\fancyhead[R]{\small\color{muted}Chance, Data, and Decisions}
\fancyfoot[L]{\footnotesize\color{muted}\edition}
\fancyfoot[R]{\footnotesize\color{muted}Page \thepage\ of \pageref*{quiz-last-page}}
\renewcommand{\headrulewidth}{.3pt}

\begin{document}

{\LARGE\bfseries\color{indigo}Conditional Probability and Bayes' Rule}\par
{\small\color{muted}50 minutes \quad $\bullet$ \quad 50 points
  \quad $\bullet$ \quad \edition}\par
\ifanswers
  \textbf{Grading note.} Subpart scores sum to the problem totals. Award
  method credit for correct reasoning after an earlier error; do not deduct
  twice for the same arithmetic mistake. Suggested working time is 46 minutes,
  leaving 4 minutes to check answers.\par
\else
  \smallskip
  \textbf{Name:}\ \rule{.49\linewidth}{.4pt}\hfill
  \textbf{Date:}\ \rule{.20\linewidth}{.4pt}\par
\fi
\textbf{Instructions.} Show your reasoning and give exact fractions.
No calculator is needed. Show the Bayes' Rule setup in Problems 3(c) and 4(d).
The suggested times leave 4 minutes to check your work.

\problem{1}{Warm-up: a finite probability space}{8}{7}
A bag contains eight tickets labeled
$a,b,b,c,c,d,d,d$. Choose one ticket uniformly at random and record its letter.
An event is any subset of the sample space of recorded letters.
\begin{enumerate}
  \item Give the sample space $\Omega$ of recorded letters and the probability
  of each outcome. Check that these probabilities sum to $1$. \partpts{4}
  \begin{solution}[.6in]
    $\Omega=\{a,b,c,d\}$, with
    $\Prob(\{a\})=1/8$, $\Prob(\{b\})=2/8$,
    $\Prob(\{c\})=2/8$, and $\Prob(\{d\})=3/8$.
    The masses are nonnegative and sum to $(1+2+2+3)/8=1$.
    For any event, add the masses of its letters; disjoint events have no
    letters in common, so their probabilities add.
    \credit{Sample space: 1; outcome probabilities: 2; sum check: 1.}
  \end{solution}

  \item Let $A=\{a,c\}$ and $B=\{c,d\}$. Find
  $\Prob(A\cap B)$ and $\Prob(A\cup B)$. \partpts{2}
  \begin{solution}[.55in]
    $A\cap B=\{c\}$, so $\Prob(A\cap B)=2/8=1/4$.
    Also $A\cup B=\{a,c,d\}$, so $\Prob(A\cup B)=(1+2+3)/8=3/4$.
    \credit{Each event probability: 1.}
  \end{solution}

  \item A student claims $\Prob(A)=2/4$ because $A$ contains two of the four
  possible letters. Is this correct? Give $\Prob(A)$ and explain. \partpts{2}
  \begin{solution}[.55in]
    No: $\Prob(A)=(1+2)/8=3/8$. The tickets are equally likely, but the
    recorded letters are not. Counting letters and dividing by four would
    require all four letters to be equally likely.
    \credit{Correct probability: 1; unequal letter probabilities: 1.}
  \end{solution}
\end{enumerate}

\problem{2}{Which group is the denominator?}{10}{10}
The table describes a group of 60 students. Choose one of them uniformly at
random. Let $C$ mean ``belongs to the chess club'' and $M$ mean ``plays a
musical instrument.''
\begin{center}
\begin{tabular}{@{}lrrr@{}}
\toprule
 & $M$ & $M^c$ & Total\\
\midrule
$C$   & 12 & 18 & 30\\
$C^c$ &  8 & 22 & 30\\
\midrule
Total & 20 & 40 & 60\\
\bottomrule
\end{tabular}
\end{center}
\begin{enumerate}
  \item Find $\Prob(C)$, $\Prob(M)$, and $\Prob(C\cap M)$. \partpts{3}
  \begin{solution}[.4in]
    $\Prob(C)=30/60=1/2$, $\Prob(M)=20/60=1/3$, and
    $\Prob(C\cap M)=12/60=1/5$.
    \credit{Each probability: 1.}
  \end{solution}

  \item Find $\Prob(M\mid C)$ and $\Prob(C\mid M)$.
  For each calculation, say which group supplies the denominator. \partpts{4}
  \begin{solution}[.65in]
    $\Prob(M\mid C)=12/30=2/5$: the denominator counts the 30 chess-club
    members. $\Prob(C\mid M)=12/20=3/5$: the denominator counts the 20
    students who play an instrument.
    \credit{Each probability: 1; each conditioning group: 1.}
  \end{solution}

  \item Are $C$ and $M$ independent? Justify your answer using probabilities.
  \partpts{3}
  \begin{solution}[.6in]
    No: $\Prob(C\cap M)=1/5$, whereas
    $\Prob(C)\Prob(M)=(1/2)(1/3)=1/6$.
    Equivalently, $\Prob(M\mid C)=2/5\ne1/3=\Prob(M)$.
    \credit{A valid independence criterion: 1; numerical comparison: 1;
    conclusion: 1. Reversing the condition and getting different values
    is not by itself an independence test.}
  \end{solution}
\end{enumerate}

\newpage
\problem{3}{A red ball: which box?}{12}{11}
Box $U$ contains three red balls and one blue ball. Box $V$ contains one red
ball and three blue balls. Choose box $U$ with probability $1/3$ and box $V$
with probability $2/3$. Then draw one ball uniformly at random from the chosen
box. Let $U$ and $V$ also denote the events that the corresponding box is
chosen, and let $R$ denote the event that the ball is red.
\begin{enumerate}
  \item List the four possible outcomes of the form (chosen box, ball color)
  and find their probabilities. Show how the probabilities for box choice and
  color combine. A labeled tree or table is fine. \partpts{4}
  \begin{solution}[2in]
    The conditional red/blue probabilities are $3/4,1/4$ in box $U$ and
    $1/4,3/4$ in box $V$. Multiply along each path:
    \begin{center}
    \begin{tabular}{@{}llc@{}}
    \toprule
    Chosen box & Color & Joint probability\\
    \midrule
    $U$ & Red  & $(1/3)(3/4)=1/4$\\[3pt]
    $U$ & Blue & $(1/3)(1/4)=1/12$\\[3pt]
    $V$ & Red  & $(2/3)(1/4)=1/6$\\[3pt]
    $V$ & Blue & $(2/3)(3/4)=1/2$\\
    \bottomrule
    \end{tabular}
    \end{center}
    These four disjoint outcomes have total probability
    $1/4+1/12+1/6+1/2=1$.
    \credit{Each outcome with its product and probability: 1.}
  \end{solution}

  \item Use the law of total probability to find $\Prob(R)$. \partpts{3}
  \begin{solution}[1.2in]
    The box choices $U,V$ are disjoint and cover all possibilities, so
    \[
      \Prob(R)=\Prob(R\mid U)\Prob(U)+\Prob(R\mid V)\Prob(V)
      =\frac34\cdot\frac13+\frac14\cdot\frac23=\frac5{12}.
    \]
    \credit{Decomposition into the two box cases: 1; correct weighted
    terms: 1; result: 1.}
  \end{solution}

  \item The drawn ball is red. Use Bayes' Rule to find the probability
  that box $U$ was chosen. \partpts{3}
  \begin{solution}[1.35in]
    \[
      \Prob(U\mid R)
      =\frac{\Prob(R\mid U)\Prob(U)}{\Prob(R)}
      =\frac{(3/4)(1/3)}{5/12}
      =\frac{1/4}{5/12}=\frac35.
    \]
    \credit{Bayes setup with the prior in the numerator: 1; correct
    substitution (or consistent use of part (b)): 1; evaluation: 1.}
  \end{solution}

  \item Another student answers part (c) with $3/4$, saying ``three of the
  four balls in box $U$ are red.'' Which conditional probability does $3/4$
  describe, and how does it differ from the probability requested? \partpts{2}
  \begin{solution}[1.15in]
    $3/4$ is $\Prob(R\mid U)$: the probability of red when the box is known
    to be $U$. The question asks for $\Prob(U\mid R)$: the probability the
    box was $U$ when the color is known to be red. The box choice is
    uncertain, and both boxes can produce a red ball.
    \credit{Identify $\Prob(R\mid U)$: 1; explain the reversed conditioning: 1.}
  \end{solution}
\end{enumerate}

\newpage
\problem{4}{Monty Hall, step by step}{20}{18}
A car is placed behind one of three doors, with probability $1/3$ for each
door. The other two doors hide goats. You initially choose Door 1.
The host follows these rules:
\begin{itemize}
  \item The host knows where the car is and always opens exactly one
  unchosen door that hides a goat.
  \item If both unchosen doors hide goats, the host chooses between them
  with equal probability $1/2$.
  \item The host always offers you the choice of staying with Door 1 or
  switching to the other unopened door.
\end{itemize}
\textbf{You observe the host open Door 3, revealing a goat.}
Let $C_i$ mean ``the car is behind Door $i$,'' and let $H_3$ mean
``the host opens Door 3.'' All probabilities below assume your initial
choice is Door 1.
\begin{enumerate}
  \item Before the host opens a door, what are
  $\Prob(C_1)$, $\Prob(C_2)$, and $\Prob(C_3)$? \partpts{2}
  \begin{solution}[.35in]
    $\Prob(C_1)=\Prob(C_2)=\Prob(C_3)=1/3$.
    \credit{Equal probabilities: 1; correct value $1/3$: 1.}
  \end{solution}

  \item Fill in the table and briefly explain each entry using the host's
  rules. These are probabilities of the host's action \emph{given the car's
  location}. \partpts{6}
  \begin{center}
  \renewcommand{\arraystretch}{1.2}
  \begin{tabular}{@{}l*{3}{>{\centering\arraybackslash}p{.95in}}@{}}
  \toprule
  Car location & $C_1$ & $C_2$ & $C_3$\\
  \midrule
  $\Prob(H_3\mid C_i)$ &
    \ifanswers $1/2$\else\rule{0pt}{.25in}\fi &
    \ifanswers $1$\fi &
    \ifanswers $0$\fi\\
  \bottomrule
  \end{tabular}
  \end{center}
  \begin{solution}[1in]
    If $C_1$ occurs, Doors 2 and 3 both hide goats, so the host chooses
    Door 3 with probability $1/2$. If $C_2$ occurs, the host must open
    Door 3. If $C_3$ occurs, opening Door 3 would reveal the car and is
    forbidden by the rules.
    \credit{Each conditional probability: 1; each explanation: 1.}
  \end{solution}

  \item Use the law of total probability to find $\Prob(H_3)$.
  Show all three car-location terms. \partpts{4}
  \begin{solution}[.95in]
    The three car locations form disjoint cases covering all possibilities:
    \[
      \Prob(H_3)=\sum_{i=1}^3\Prob(H_3\mid C_i)\Prob(C_i)
      =\frac12\cdot\frac13+1\cdot\frac13+0\cdot\frac13=\frac12.
    \]
    \credit{Total-probability setup: 1; three weighted terms: 2;
    evaluation: 1.}
  \end{solution}

  \item Use Bayes' Rule to find $\Prob(C_1\mid H_3)$ and
  $\Prob(C_2\mid H_3)$. Also state $\Prob(C_3\mid H_3)$. \partpts{5}
  \begin{solution}[1.25in]
    \[
      \Prob(C_i\mid H_3)
      =\frac{\Prob(H_3\mid C_i)\Prob(C_i)}{\Prob(H_3)}.
    \]
    Thus
    \[
      \Prob(C_1\mid H_3)=\frac{(1/2)(1/3)}{1/2}=\frac13,
      \qquad
      \Prob(C_2\mid H_3)=\frac{1(1/3)}{1/2}=\frac23.
    \]
    $\Prob(C_3\mid H_3)=0$, since Door 3 has revealed a goat.
    The three updated probabilities sum to $1$.
    \credit{For each of Doors 1 and 2, Bayes setup: 1 and calculation: 1;
    Door 3 probability: 1.}
  \end{solution}

  \item Should you stay with Door 1 or switch to Door 2? Give the chance
  of winning under each choice. Explain why ``two doors remain, so each
  has probability $1/2$'' is not a valid argument. \partpts{3}
  \begin{solution}[1in]
    Switch to Door 2: switching wins with probability $2/3$, while staying
    wins with probability $1/3$, given $H_3$. Two remaining possibilities
    need not be equally likely. Opening Door 3 is twice as likely when
    the car is behind Door 2 as when it is behind Door 1; with equal priors,
    Bayes' Rule therefore gives Door 2 twice the updated probability.
    \credit{Switch decision: 1; both winning probabilities: 1; explanation
    referring to unequal likelihoods or the host's informed rule: 1.}
  \end{solution}
\end{enumerate}

\label{quiz-last-page}
\end{document}
